\begin{remarks}[5.6.9]
\label{IV.2.s5.6.9-fr}
\begin{enumerate}
\item[(i)] Incluso si $Y$ es regular, $X$ irreducible y $f:X\to Y$ dominante
y de tipo finito, los dos miembros de (5.6.6.1) no son necesariamente iguales,
como muestra el ejemplo en que $Y=\operatorname{Spec}(A)$, donde $A$ es un
anillo de valoración discreta, $X=\operatorname{Spec}(K)$, donde $K$ es el
cuerpo de fracciones de $A$, y $f$ es el morfismo canónico.
\item[(ii)] El ejemplo (5.4.3, (i)) muestra que en (5.6.6) y (5.6.8) no se
puede suprimir la hipótesis de que $X$ es irreducible, verificándose las demás
hipótesis. Veremos, sin embargo, (10.6.1) que, bajo hipótesis adicionales sobre
$Y$, satisfechas por ejemplo cuando $Y$ es un preesquema localmente de tipo
finito sobre un cuerpo, o sobre un anillo de Dedekind que tiene una infinidad
de ideales primos, tales fenómenos no pueden presentarse.
\end{enumerate}
\end{remarks}

\begin{proposition}[5.6.10]
\label{IV.2.s5.6.10-fr}
Sean $A$ un anillo local noetheriano íntegro universalmente catenario y $B$
un anillo íntegro que contiene a $A$ y que es una $A$-álgebra finita. Entonces,
para todo ideal maximal $\mathfrak n$ de $B$, se tiene
$\dim(B_{\mathfrak n})=\dim(A)$.
\end{proposition}

En efecto, se tiene $\operatorname{deg.tr}_A B=0$ y el cuerpo residual $k'$
de $B_{\mathfrak n}$ es una extensión algebraica del cuerpo de fracciones de
$A$. La conclusión se sigue de la fórmula (5.6.1.1), pues todo ideal maximal
de $B$ está situado sobre el de $A$.

